\section{Try It Yourself}

Reading builds understanding; breaking things builds intuition. Each exercise
takes ten to thirty minutes.

\begin{enumerate}
\item \textbf{Watch a lease expire.} Start \texttt{make dev}, submit a job with
  \texttt{python3 scripts/smoke.py}, and while it runs kill the local agent
  (\texttt{pkill -f agent.py}). Query the job every few seconds:
  \texttt{curl -H "X-Dev-User: alice@uni.edu" localhost:8080/jobs/<id>}. Watch
  the attempt go \texttt{running}, \texttt{unknown}, \texttt{failed
  (lease\_lost)}, and a second attempt start.

\item \textbf{Be a zombie.} In a Python shell, use \texttt{harness.client.FakeAgent}
  to claim a job, let its lease expire, let a second fake agent complete it,
  then call \texttt{complete} with the first agent's token. Read the
  \texttt{409} body.

\item \textbf{See SKIP LOCKED.} Open two \texttt{psql} sessions. In the first,
  \texttt{BEGIN; SELECT id FROM jobs ... FOR UPDATE SKIP LOCKED LIMIT 1;} and
  do not commit. In the second run the same query. Note which row each gets.
  Now remove \texttt{SKIP LOCKED} from the second and see it wait.

\item \textbf{Stream logs by hand.}
  \texttt{curl -N -H "X-Dev-User: alice@uni.edu"
  localhost:8080/jobs/<id>/logs/stream} shows raw SSE. Stop it, then rerun with
  \texttt{-H "Last-Event-ID: 3"}.

\item \textbf{Verify a signature by hand.} Fetch
  \texttt{/agent/bundle-key} and \texttt{/agent/bundle} (with a node token),
  then call \texttt{agent.updater.verify\_signature}. Flip one byte of the bundle
  and verify again.

\item \textbf{Break an update on purpose.} In the simulated lab, add
  \texttt{raise SystemExit(1)} at the top of \texttt{agent/agent.py}, bump the
  version, rebuild only \texttt{sim-control}, and watch the PCs update, fail
  three times, roll back, and list the version in
  \texttt{/opt/lab-agent/rejected}. Then revert.

\item \textbf{Add a scenario to the exam.} Write a \texttt{sim/exam.py}
  scenario that pauses a PC (\texttt{docker pause}) for longer than a lease and
  checks what happens to its job. What should happen when it resumes?

\item \textbf{Add a migration.} Add a nullable \texttt{notes} column to
  \texttt{nodes}, autogenerate a migration, read it, run \texttt{make test}, and
  confirm \texttt{alembic check} is clean.
\end{enumerate}

\subsection{Where these ideas show up next}

\begin{description}
\item[HTCondor] a decades-old system built for exactly this, scavenging idle
  desktops for computation, with policies that respect the desktop's owner.
\item[Slurm] the standard scheduler on university supercomputers; partitions,
  QOS limits and GRES map to this project's classes and quotas.
\item[Kubernetes] node labels, resource requests, leases for leader election,
  rolling updates and readiness checks.
\item[CI systems] GitHub Actions and GitLab runners pull jobs, run them in
  containers, stream logs and upload artifacts, the same loop as this agent.
\item[Render farms and BOINC] distributing independent work units to many
  machines and validating the results.
\item[Payment APIs] idempotency keys and exactly-once effects.
\end{description}

\subsection{Further reading}

\begin{itemize}
\item Martin Kleppmann, \emph{Designing Data-Intensive Applications}, chapters
  on replication, transactions and ``the trouble with distributed systems''.
\item Martin Kleppmann, ``How to do distributed locking'' (fencing tokens).
\item The PostgreSQL manual, ``Explicit Locking'' and the \texttt{SELECT}
  locking clauses.
\item The Twelve-Factor App (\texttt{12factor.net}).
\item Google, \emph{Site Reliability Engineering}, on alerting, runbooks and
  release engineering.
\end{itemize}
